\documentclass[11pt,letterpaper,oneside]{report}
\usepackage[T1]{fontenc}
\usepackage[utf8]{inputenc}
\usepackage{lmodern}
\usepackage[margin=0.90in,headheight=15pt]{geometry}
\usepackage{microtype}
\usepackage[svgnames,table]{xcolor}
\usepackage{amsmath,amssymb,booktabs,longtable,tabularx,array}
\usepackage{enumitem}
\usepackage[most]{tcolorbox}
\usepackage{listings}
\usepackage{fancyhdr}
\usepackage{titlesec}
\usepackage{hyperref}
\usepackage{xurl}
\usepackage{lastpage}
\definecolor{ink}{HTML}{14273C}
\definecolor{primary}{HTML}{0B5379}
\definecolor{accent}{HTML}{B66A30}
\definecolor{pale}{HTML}{F1F5F8}
\definecolor{codebg}{HTML}{F5F7F8}
\hypersetup{colorlinks=true,linkcolor=primary,urlcolor=primary,pdftitle={Linux and UNIX System Programming Engineering Suite}}
\titleformat{\chapter}[display]{\normalfont\LARGE\bfseries\color{ink}}{\color{accent}\large\MakeUppercase{\chaptername}\ \thechapter}{10pt}{}[\vspace{0.4ex}\titlerule]
\titleformat{\section}{\Large\bfseries\color{primary}}{\thesection}{0.7em}{}
\titleformat{\subsection}{\large\bfseries\color{ink}}{\thesubsection}{0.7em}{}
\setlist[itemize]{leftmargin=*,topsep=3pt,itemsep=2pt}
\setlist[enumerate]{leftmargin=*,topsep=3pt,itemsep=3pt}
\setlength{\parskip}{5pt}
\setlength{\parindent}{0pt}
\lstdefinestyle{sysc}{language=C,basicstyle=\ttfamily\footnotesize,backgroundcolor=\color{codebg},keywordstyle=\bfseries\color{primary},commentstyle=\color{ForestGreen},stringstyle=\color{Brown},breaklines=true,showstringspaces=false,tabsize=4,frame=single,rulecolor=\color{LightSteelBlue},columns=fullflexible,keepspaces=true}
\lstset{style=sysc}
\newtcolorbox{engineeringnote}[1][]{colback=pale,colframe=primary,boxrule=0.55pt,arc=1.5mm,title=Engineering note,fonttitle=\bfseries,#1}
\newtcolorbox{modernnote}[1][]{colback=white,colframe=accent,boxrule=0.55pt,arc=1.5mm,title=Modern extension --- not in the 2010 textbook,fonttitle=\bfseries,#1}
\newcommand{\sys}[1]{\texttt{#1}}
\pagestyle{fancy}\fancyhf{}\fancyhead[L]{\small\color{ink} Linux / UNIX Engineering Suite}\fancyhead[R]{\small\color{primary}\nouppercase{\leftmark}}\fancyfoot[C]{\small\thepage\ / \pageref{LastPage}}
\renewcommand{\headrulewidth}{0.45pt}
\setcounter{tocdepth}{1}
\setlength{\emergencystretch}{2em}
\begin{document}
\begin{titlepage}\centering\vspace*{2cm}{\color{primary}\Large ENGINEERING REFERENCE SERIES\par}\vspace{1.25cm}{\color{ink}\Huge\bfseries C / POSIX Implementation Handbook\par}\vspace{0.6cm}{\Large Reference patterns with complete reproducible programs\par}\vspace{2cm}\rule{0.75\textwidth}{1.2pt}\par\vspace{0.55cm}{\large An original companion to Michael Kerrisk\textquotesingle s 2010 text}\par\vspace{0.3cm}{\large Edition 1.0\quad |\quad October 2026}\par\vfill{\small Document 03\quad |\quad pdfLaTeX-compatible standalone source}\par\end{titlepage}
\tableofcontents\clearpage
\chapter{C interfaces and build discipline}
The programs in this handbook are original minimal examples that make important invariants visible; they are not production-ready libraries. They assume a GCC/Clang toolchain on a recent Linux system, and explicitly identify Linux-specific interfaces. The source archive also contains the same listings under examples/ for compilation and experimentation.

\section{Compiler, warning and instrumentation baseline}
\begin{lstlisting}
cc -std=c11 -O2 -g -Wall -Wextra -Wpedantic -Werror \
   -D_FORTIFY_SOURCE=2 examples/checked_copy.c -o checked_copy
cc -std=c11 -O1 -g -Wall -Wextra -Wpedantic \
   -fsanitize=address,undefined examples/spawn_wait.c -o spawn_wait
cc -std=c11 -O2 -g -Wall -Wextra -Wpedantic \
   -pthread examples/pthread_bounded_queue.c -o pthread_queue
\end{lstlisting}
\section{Error-return taxonomy}
\begin{itemize}
\item Most syscall wrappers: -1 and errno set; callers should capture errno before invoking diagnostic functions that may modify it.
\item POSIX pthread APIs: usually return an error number directly rather than setting errno; report rc, not a stale errno.
\item getaddrinfo: returns EAI\_* status; use gai\_strerror for messages.
\item Some APIs allow partial progress before an error. Never discard bytes already transferred or repeat a non-idempotent operation blindly.
\end{itemize}
\chapter{Binary-safe I/O and cleanup}
This copier creates a destination exclusively, defaults its mode to 0600 and correctly handles short writes and EINTR. It is not atomic file replacement; successful fsync of the file is not, by itself, a complete power-loss durability protocol for directory entries.

\begin{lstlisting}
#define _POSIX_C_SOURCE 200809L
#include <errno.h>
#include <fcntl.h>
#include <stdio.h>
#include <stdlib.h>
#include <unistd.h>

static int copy_stream(int src, int dst) {
    unsigned char buf[16384];
    for (;;) {
        ssize_t n = read(src, buf, sizeof buf);
        if (n < 0 && errno == EINTR) continue;
        if (n < 0) return -1;
        if (n == 0) return 0;
        size_t used = 0;
        while (used < (size_t)n) {
            ssize_t w = write(dst, buf + used, (size_t)n - used);
            if (w < 0 && errno == EINTR) continue;
            if (w <= 0) { if (w == 0) errno = EIO; return -1; }
            used += (size_t)w;
        }
    }
}
int main(int argc, char **argv) {
    if (argc != 3) {
        fprintf(stderr, "usage: %s SOURCE NEW-DESTINATION\n", argv[0]);
        return EXIT_FAILURE;
    }
    int in = open(argv[1], O_RDONLY | O_CLOEXEC);
    if (in < 0) { perror("open source"); return EXIT_FAILURE; }
    int out = open(argv[2], O_WRONLY | O_CREAT | O_EXCL | O_CLOEXEC,
                   0600);
    if (out < 0) { perror("open destination"); close(in); return EXIT_FAILURE; }
    int ok = copy_stream(in, out);
    int saved = errno;
    if (ok == 0 && fsync(out) < 0) { ok = -1; saved = errno; }
    if (close(out) < 0 && ok == 0) { ok = -1; saved = errno; }
    if (close(in) < 0 && ok == 0) { ok = -1; saved = errno; }
    if (ok < 0) {
        errno = saved; perror("copy");
        if (unlink(argv[2]) < 0) perror("cleanup destination");
        return EXIT_FAILURE;
    }
    return EXIT_SUCCESS;
}
\end{lstlisting}
\section{Test oracle}
\begin{lstlisting}
printf 'abc\000def' > sample.bin
./checked_copy sample.bin copied.bin
cmp sample.bin copied.bin
./checked_copy sample.bin copied.bin  # must fail: exclusive creation
\end{lstlisting}
\begin{engineeringnote}
A created destination is removed on failure. A production transactional copy may require writing to a unique temporary file within a trusted destination directory, syncing file contents, renaming atomically, and syncing the containing directory with filesystem-specific assumptions documented.
\end{engineeringnote}
\chapter{Controlled child launch and reaping}
Avoid shell interpolation when arguments are structured data. A fixed executable path and deliberately constructed environment reduce unexpected execution dependencies. The sample checks both the spawn return code and child exit status.

\begin{lstlisting}
#define _POSIX_C_SOURCE 200809L
#include <errno.h>
#include <spawn.h>
#include <stdio.h>
#include <stdlib.h>
#include <sys/wait.h>
#include <unistd.h>

int main(void) {
    pid_t pid;
    char *const args[] = { "printf", "worker-ready\n", NULL };
    char *const env[] = { "LC_ALL=C", "PATH=/usr/bin:/bin", NULL };
    int rc = posix_spawn(&pid, "/usr/bin/printf", NULL, NULL, args, env);
    if (rc != 0) { errno = rc; perror("posix_spawn"); return 1; }
    int status;
    for (;;) {
        if (waitpid(pid, &status, 0) >= 0) break;
        if (errno == EINTR) continue;
        perror("waitpid"); return 1;
    }
    if (!WIFEXITED(status) || WEXITSTATUS(status) != 0) {
        fprintf(stderr, "child failed\n"); return 1;
    }
    return 0;
}
\end{lstlisting}
\section{Review the contract}
\begin{itemize}
\item What file descriptors should survive exec?
\item How are executable absence and permission failure reported?
\item What is the required behavior when the parent receives a signal while waiting?
\item How would you communicate timeouts without mixing asynchronous signal handlers into all business logic?
\end{itemize}
\chapter{Shared-memory concurrency with pthreads}
Condition variables wait for state changes, not for a reliable count of notifications. Always test the predicate in a loop while the same mutex protects the queue. This example uses a single consumer to make sum validation deterministic; extending it to multiple producers and consumers requires an explicit close/wakeup protocol.

\begin{lstlisting}
#define _POSIX_C_SOURCE 200809L
#include <pthread.h>
#include <stdio.h>
#include <stdlib.h>

#define N 8
struct queue {
    int a[N], head, tail, count, closed;
    pthread_mutex_t mu;
    pthread_cond_t readable, writable;
};
static struct queue q = { .mu = PTHREAD_MUTEX_INITIALIZER,
    .readable = PTHREAD_COND_INITIALIZER,
    .writable = PTHREAD_COND_INITIALIZER };
static void *consumer(void *unused) {
    (void)unused;
    long sum = 0;
    pthread_mutex_lock(&q.mu);
    for (;;) {
        while (q.count == 0 && !q.closed)
            pthread_cond_wait(&q.readable, &q.mu);
        if (q.count == 0 && q.closed) break;
        int x = q.a[q.head]; q.head = (q.head + 1) % N; --q.count;
        pthread_cond_signal(&q.writable);
        pthread_mutex_unlock(&q.mu);
        sum += x;
        pthread_mutex_lock(&q.mu);
    }
    pthread_mutex_unlock(&q.mu);
    printf("sum=%ld\n", sum);
    return NULL;
}
int main(void) {
    pthread_t t;
    int rc = pthread_create(&t, NULL, consumer, NULL);
    if (rc != 0) { fprintf(stderr, "pthread_create: %d\n", rc); return 1; }
    for (int i = 1; i <= 100; ++i) {
        pthread_mutex_lock(&q.mu);
        while (q.count == N)
            pthread_cond_wait(&q.writable, &q.mu);
        q.a[q.tail] = i; q.tail = (q.tail + 1) % N; ++q.count;
        pthread_cond_signal(&q.readable);
        pthread_mutex_unlock(&q.mu);
    }
    pthread_mutex_lock(&q.mu);
    q.closed = 1;
    pthread_cond_broadcast(&q.readable);
    pthread_mutex_unlock(&q.mu);
    rc = pthread_join(t, NULL);
    if (rc != 0) { fprintf(stderr, "pthread_join: %d\n", rc); return 1; }
    pthread_cond_destroy(&q.writable); pthread_cond_destroy(&q.readable);
    pthread_mutex_destroy(&q.mu);
    return 0;
}
\end{lstlisting}
\section{Correctness checks}
\begin{lstlisting}
./pthread_queue     # must output sum=5050
# Run in a loop under the race/undefined-behavior tools available.
# Check all queue operations hold q.mu and all waits test predicates.
\end{lstlisting}
\chapter{Local sockets}
socketpair creates connected UNIX-domain endpoints without a pathname. It is useful for process-local testing, but a single read from a stream socket is not a general message-framing protocol.

\begin{lstlisting}
#define _POSIX_C_SOURCE 200809L
#include <errno.h>
#include <stddef.h>
#include <stdio.h>
#include <stdlib.h>
#include <string.h>
#include <sys/socket.h>
#include <sys/un.h>
#include <unistd.h>

int main(void) {
    int fds[2];
    if (socketpair(AF_UNIX, SOCK_STREAM, 0, fds) == -1) {
        perror("socketpair"); return 1;
    }
    const char msg[] = "bounded local message";
    if (write(fds[0], msg, sizeof msg) != (ssize_t)sizeof msg) {
        perror("write"); close(fds[0]); close(fds[1]); return 1;
    }
    char buf[128];
    ssize_t n = read(fds[1], buf, sizeof buf);
    if (n < 0) perror("read");
    else printf("received %zd bytes\n", n);
    close(fds[0]); close(fds[1]);
    return n < 0;
}
\end{lstlisting}
\chapter{Linux readiness notifications}
eventfd and epoll are Linux-specific. This minimal event-loop experiment writes a counter and waits for readiness, then consumes the counter.

\begin{lstlisting}
#define _GNU_SOURCE
#include <errno.h>
#include <stdint.h>
#include <stdio.h>
#include <stdlib.h>
#include <sys/epoll.h>
#include <sys/eventfd.h>
#include <unistd.h>

int main(void) {
    int e = eventfd(0, EFD_CLOEXEC | EFD_NONBLOCK);
    if (e == -1) { perror("eventfd"); return 1; }
    int ep = epoll_create1(EPOLL_CLOEXEC);
    if (ep == -1) { perror("epoll_create1"); close(e); return 1; }
    struct epoll_event event = { .events = EPOLLIN,
                                 .data.fd = e };
    if (epoll_ctl(ep, EPOLL_CTL_ADD, e, &event) == -1) {
        perror("epoll_ctl"); close(ep); close(e); return 1;
    }
    uint64_t one = 1, got = 0;
    if (write(e, &one, sizeof one) != (ssize_t)sizeof one) {
        perror("write eventfd"); close(ep); close(e); return 1;
    }
    struct epoll_event ready;
    int count = epoll_wait(ep, &ready, 1, 1000);
    if (count != 1 || !(ready.events & EPOLLIN) ||
        read(e, &got, sizeof got) != (ssize_t)sizeof got) {
        fprintf(stderr, "readiness test failed\n");
        close(ep); close(e); return 1;
    }
    printf("event count=%llu\n", (unsigned long long)got);
    close(ep); close(e); return 0;
}
\end{lstlisting}
\chapter{Fd-anchored file inspection}
The example uses an already-opened directory as a trust anchor and refuses a symbolic link for the leaf. It does not claim comprehensive prevention of traversal through attacker-controlled intermediate directories; Linux openat2 can provide more expressive resolution policies when supported.

\begin{lstlisting}
#define _POSIX_C_SOURCE 200809L
#include <errno.h>
#include <fcntl.h>
#include <stdio.h>
#include <stdlib.h>
#include <sys/stat.h>
#include <unistd.h>

/* Caller must already hold a trustworthy directory fd.
   This prevents symlinks at the final component, but not all
   possible traversal across attacker-controlled parents. */
static int read_regular_leaf(int dirfd, const char *leaf) {
    int fd = openat(dirfd, leaf, O_RDONLY | O_CLOEXEC | O_NOFOLLOW);
    if (fd < 0) return -1;
    struct stat st;
    if (fstat(fd, &st) < 0) {
        int saved = errno;
        close(fd);
        errno = saved;
        return -1;
    }
    if (!S_ISREG(st.st_mode)) {
        close(fd);
        errno = EINVAL;
        return -1;
    }
    char c;
    ssize_t n = read(fd, &c, 1);
    int saved = errno;
    close(fd);
    errno = saved;
    return n >= 0 ? 0 : -1;
}
int main(void) {
    int dir = open(".", O_RDONLY | O_DIRECTORY | O_CLOEXEC);
    if (dir < 0) { perror("open dir"); return 1; }
    int rc = read_regular_leaf(dir, "input.txt");
    if (rc < 0) perror("read_regular_leaf");
    close(dir); return rc < 0;
}
\end{lstlisting}
\chapter{Productionization criteria}
\begin{enumerate}
\item Separate reusable components from main and add narrow unit-testable APIs.
\item Document all ownership, allowed concurrency and nonblocking behavior.
\item Apply explicit size, count, path and time limits.
\item Introduce fault injection (short writes, EINTR, EAGAIN, absent resources, peer closure).
\item Run sanitizer, static analysis and integration coverage, including negative-path cleanup.
\item Record portability and distribution requirements for each API.
\end{enumerate}

\chapter*{Source and currency notes}
\addcontentsline{toc}{chapter}{Source and currency notes}
\textbf{Foundational source:} Michael Kerrisk, \emph{The Linux Programming Interface: A Linux and UNIX System Programming Handbook}, No Starch Press, 2010, ISBN 978-1-59327-220-3. The textbook's chapter titles and start pages organize the study guide; explanations, tasks, and code in this suite are independently authored educational material, not reproductions of the original text or its example programs.

\textbf{Current reference baseline:} The Open Group / IEEE, \emph{POSIX.1-2024, Issue 8}, \url{https://pubs.opengroup.org/onlinepubs/9799919799/}; Linux man-pages project, \url{https://man7.org/linux/man-pages/}; Linux kernel documentation, \url{https://docs.kernel.org/}. For Linux-only extensions consult the relevant per-interface manual pages and kernel documentation for the target kernel, architecture and libc; availability can differ materially.

\textbf{Scope:} This is an engineering study and implementation companion. It does not claim to restate a normative standard or guarantee a specific Linux kernel's behavior. System calls, libc wrappers and higher-level library functions are explicitly distinguished where material.
\end{document}
